\documentclass[11pt]{article}
\usepackage[margin=1in]{geometry}
\usepackage{amsmath,amssymb,amsthm}
\usepackage{hyperref}
\usepackage{xurl}

\newtheorem{theorem}{Theorem}
\newtheorem{lemma}[theorem]{Lemma}
\newtheorem{corollary}[theorem]{Corollary}
\theoremstyle{definition}
\newtheorem{definition}[theorem]{Definition}
\theoremstyle{remark}
\newtheorem{remark}[theorem]{Remark}

\let\oldus\_
\renewcommand{\_}{\oldus\allowbreak}
\newcommand{\R}{\mathbb{R}}
\newcommand{\Hm}{\mathcal{H}^1}

\title{A proof of Conjecture 00000006420\\[2pt]
\large Two unions of two unit disks with the same isoperimetric deficit and different Fraenkel asymmetry}
\author{Jackmeson1}
\date{}

\begin{document}
\sloppy
\maketitle

\section{The conjecture}

Conjecture 00000006420 reads (English version; the Chinese version agrees):

\begin{quote}
\emph{Definition:} The isoperimetric deficit and asymmetry are two layers.
\emph{Conjecture:} There exist two sets with the same deficit but different asymmetry, and the
separation is realized by an explicit two-circle union with the same deficit but different
asymmetry.
\end{quote}

This is an existence claim. We exhibit two explicit sets in the plane, each the union of two
disjoint closed unit disks, prove that they have the same isoperimetric deficit (namely
$\sqrt2-1$), and prove that their Fraenkel asymmetries differ. The main results are formalized in
Lean 4 / Mathlib for the perimeter $P(E)=\mathcal H^1(\partial E)$ defined below. The De Giorgi
(distributional) perimeter is not defined in Lean: \texttt{perimeter\_axiomatic} is a conditional
statement about any translation-invariant perimeter that is additive on disjoint compact sets, and
its application to the De Giorgi perimeter (and that perimeter's agreement with $\mathcal H^1(\partial
E)$ for these sets) is a standard fact that is not formalized; see the conventions paragraph at the
end.

\section{Reading and definitions}

We work in the Euclidean plane $\R^2$ (Mathlib: \texttt{EuclideanSpace R (Fin 2)}) with Lebesgue
measure $|\cdot|$. ``Asymmetry'' is read as the Fraenkel asymmetry, the asymmetry index that is
paired with the isoperimetric deficit in the quantitative isoperimetric inequality
$\alpha(E)^2\le C\,\delta(E)$. We use the normalizations of Cicalese and Leonardi
(\emph{A selection principle for the sharp quantitative isoperimetric inequality},
\url{https://arxiv.org/abs/1007.3899}, Section 2), which follow Fusco, Maggi and Pratelli. A
``two-circle union'' is read as the union of two closed disks. Two circles as curves have area
zero, so their deficit is not defined.

\begin{definition}
For $E\subset\R^2$ let $r_E=\sqrt{|E|/\pi}$, the radius of the disk $B_E$ with $|B_E|=|E|$.
\begin{itemize}
\item \emph{Perimeter:} $P(E)=\Hm(\partial E)$, the one-dimensional Hausdorff measure (length) of
  the topological boundary. In Lean this is \texttt{perimeter E := muH[1] (frontier E)}. Mathlib's
  $\mu_H[1]$ needs no normalizing constant in dimension one; on $\R$ it equals Lebesgue measure
  (\texttt{hausdorffMeasure\_real}).
\item \emph{Isoperimetric deficit:} $\delta(E)=\dfrac{P(E)-P(B_E)}{P(B_E)}=\dfrac{P(E)}{P(B_E)}-1$,
  with $B_E=\bar B(0,r_E)$ (\texttt{deficit}).
\item \emph{Fraenkel asymmetry:}
  $\alpha(E)=\inf_{x\in\R^2}\dfrac{|E\,\triangle\,\bar B(x,r_E)|}{|E|}$ (\texttt{asymmetry}).
\end{itemize}
\end{definition}

\section{The two sets}

Let $D_a=\bar B((a,0),1)$ be the closed unit disk centred at $(a,0)$, and put
\[
E_{\rm near}=D_0\cup D_{5/2},\qquad E_{\rm far}=D_0\cup D_{10}.
\]
In both unions the centres are more than $2$ apart, so the two disks are disjoint
(\texttt{disj}). Hence $|E_{\rm near}|=|E_{\rm far}|=2\pi$, $r_E=\sqrt2$ in both cases, and every
comparison disk $\bar B(x,\sqrt2)$ has area $2\pi$ (\texttt{area\_near}, \texttt{area\_far},
\texttt{eqRadius\_near}, \texttt{eqRadius\_far}, \texttt{vol\_comparison}).

\section{Same deficit}

\begin{lemma}[\texttt{frontier\_union\_of\_closed}]
If $A,B$ are closed and disjoint, then $\partial(A\cup B)=\partial A\cup\partial B$.
\end{lemma}
\begin{proof}
The inclusion $\subseteq$ holds in general. Conversely, let $p\in\partial A$. Then $p\in A$, so
$p\in\overline{A\cup B}$. If $p$ were interior to $A\cup B$, then $U=\operatorname{int}(A\cup B)
\setminus B$ would be an open neighbourhood of $p$ (because $p\notin B$ and $B$ is closed)
contained in $A$, contradicting $p\in\partial A$.
\end{proof}

Let $L=\Hm(S^1)$ be the length of the unit circle.

\begin{lemma}[\texttt{perimeter\_two}, \texttt{L\_pos}, \texttt{L\_lt\_top}]
For $|a|>2$ we have $P(D_0\cup D_a)=2L$, and $0<L<\infty$. The proofs establish $L\le 2\pi$ and
$2\le L$ as intermediate steps.
\end{lemma}
\begin{proof}
By the previous lemma and $\partial \bar B(c,1)=S(c,1)$, the boundary is the disjoint union of two
unit circles. $\Hm$ is a Borel measure and is invariant under translations (isometries), so
$P=2L$. For the upper bound, every point of the unit circle in $\mathbb{C}\cong\R^2$ is
$e^{i\theta}$ with $\theta=\arg z\in[-\pi,\pi]$, and $\theta\mapsto e^{i\theta}$ is
$1$-Lipschitz, because $|e^{is}-e^{it}|=|e^{i(s-t)}-1|\le|s-t|$. Lipschitz maps do not increase
$\Hm$, so $L\le\Hm([-\pi,\pi])=2\pi$. For the lower bound, the projection $z\mapsto\operatorname{Re}z$
is $1$-Lipschitz and maps the circle onto $[-1,1]$, so $2=\Hm([-1,1])\le L$. The identification
$\mathbb{C}\cong\R^2$ is the isometry given by the orthonormal basis $(1,i)$
(\texttt{L\_eq\_complex}).
\end{proof}

\begin{theorem}[\texttt{deficit\_two}]
$\delta(E_{\rm near})=\delta(E_{\rm far})=\sqrt2-1$.
\end{theorem}
\begin{proof}
We have $\partial\bar B(0,\sqrt2)=S(0,\sqrt2)=\sqrt2\cdot S(0,1)$, and $\Hm$ scales by the factor
$\sqrt2$ under this dilation (\texttt{Measure.hausdorffMeasure\_smul0}). So
$P(B_E)=\sqrt2\,L$ (\texttt{perimeter\_comparison}), and
$\delta=2L/(\sqrt2 L)-1=\sqrt2-1$, using $0<L<\infty$.
\end{proof}

\section{Different asymmetry}

\begin{theorem}[\texttt{asymmetry\_far}]
$\alpha(E_{\rm far})\ge 1$.
\end{theorem}
\begin{proof}
Fix $x$ and let $B=\bar B(x,\sqrt2)$. The disk $B$ cannot meet both $D_0$ and $D_{10}$: otherwise
$10=|(0,0)-(10,0)|\le 1+2\sqrt2+1<5$ (\texttt{far\_cases}). Say $B\cap D_{10}=\emptyset$; the other
case is symmetric. Then $D_{10}\subseteq E\setminus B$ and $B\setminus D_0\subseteq B\setminus E$,
and these two sets are disjoint. So
$|E\triangle B|\ge |D_{10}|+|B|-|D_0|=\pi+2\pi-\pi=2\pi$ (\texttt{symmDiff\_lower}). Dividing by
$|E|=2\pi$ gives a quotient of at least $1$ for every $x$, so the infimum is at least $1$.
\end{proof}

\begin{theorem}[\texttt{near\_value}, \texttt{asymmetry\_near\_le}, \texttt{asymmetry\_near}]
$\alpha(E_{\rm near})\le 1-\tfrac1{100}<1$.
\end{theorem}
\begin{proof}
Take $x_0=(2/5,0)$, $B_0=\bar B(x_0,\sqrt2)$ and $S=\bar B((8/5,0),1/10)$. Then:
$D_0\subseteq B_0$ since $1+\tfrac25=\tfrac75\le\sqrt2$, because $\tfrac{49}{25}\le2$;
$S\subseteq D_{5/2}$ since $\tfrac1{10}+\tfrac9{10}=1$;
and $S\subseteq B_0$ since $\tfrac1{10}+\tfrac65=\tfrac{13}{10}\le\sqrt2$.
So $K=D_0\cup S\subseteq E_{\rm near}\cap B_0$ is a disjoint union with
$|K|=\pi+\pi/100$. Since $E\triangle B_0\subseteq (E\setminus K)\cup(B_0\setminus K)$, we get
\[
|E_{\rm near}\triangle B_0|\le (2\pi-|K|)+(2\pi-|K|)=2\pi-\tfrac{\pi}{50}.
\]
The quotient at $x_0$ is therefore at most $1-\tfrac1{100}$. The quotients are non-negative, so the
infimum is at most this value.
\end{proof}

\begin{theorem}[\texttt{separation}, \texttt{conjecture6420}]
$E_{\rm near}$ and $E_{\rm far}$ are unions of two disjoint closed unit disks. They have equal
area $2\pi$, equal perimeter $2L\in(0,\infty)$, and equal deficit $\sqrt2-1$, while
$\alpha(E_{\rm near})<1\le\alpha(E_{\rm far})$. Hence the conjecture holds.
\end{theorem}

\paragraph{Other conventions (\texttt{normalisations}, \texttt{perimeter\_axiomatic}).}
(i) Any deficit that is a function of $(P(E),|E|)$, for example $P(E)^2/(4\pi|E|)-1$ or
$P(E)-2\sqrt{\pi|E|}$, agrees on the two sets. (ii) Any positive rescaling of $\alpha$ still
separates them. This covers the normalization $|E\triangle(x+B_E)|/r_E^{2}$, which is $\pi$ times
ours. (iii) The De Giorgi (distributional) perimeter is the convention of the cited paper. Lean
proves that \emph{every} perimeter functional that is translation invariant and additive on
disjoint compact sets, a pair of properties that the De Giorgi perimeter has, takes the same value
on $E_{\rm near}$ and $E_{\rm far}$. So every deficit built from such a perimeter and the area
agrees. That the De Giorgi perimeter has these two properties, and that it equals $\Hm(\partial E)$
for these smooth sets, are standard facts that are not formalized here.

\section{Lean formalization and axiom audit}

The file \texttt{lean/Conjecture6420/Basic.lean} (Lean 4.33.1, Mathlib v4.33.1, single import
\texttt{Mathlib}) contains all definitions and proofs, in namespace \texttt{C6420}. Main
theorems: \texttt{separation}, \texttt{conjecture6420}, \texttt{normalisations} and
\texttt{perimeter\_axiomatic}. For each of them, \texttt{\#print axioms} reports exactly
\texttt{[propext, Classical.choice, Quot.sound]}. There is no \texttt{sorry}, no
\texttt{native\_decide} and no added axiom. The value of $L$ ($=2\pi$) is never needed: only
$0<L<\infty$ is used.

\end{document}
